\exnumber{05}
\extitle{UX II - Ex 6 - User flow}
\exfiles{ex05/zip archive, named Name\_FirstName\_DeliverableName\_Date\_VersionNumber}
\exauthorize{Those named in the exercise}
\makeheaderfiles

\textbf{Instructions:}

Structure your main ideas into a comprehensive user flow.
You can include one, two, of the three main ideas from your ideation process, depending on how consistent it would be for the users.

\textbf{Deliverable:}

\begin{itemize}
\item A comprehensive user flow on a board (digital or pictures of a paper version)
\item Upload the documents on Github, in a zip folder named "ProjectTitle\_Name\_FirstName", with all project deliverables as follows: "Name\_FirstName\_DeliverableName\_Date\_VersionNumber".  For example, the first deliverable can be named as "Doe\_John\_MidFiPrototype\_05242024\_V1".
\item Add a title within all written documents.
\item Pay attention to the global aspect of the deliverables. They must be clear, simple and easy to read. You can get inspiration on canva.com but don't overload your design with too much details!
\end{itemize}

\textbf{Resources:}

\begin{itemize}
\item User flow glossary \url{https://www.productplan.com/glossary/user-flow/}
\item Site flows vs User Flows \url{https://uxmovement.com/wireframes/site-flows-vs-user-flows-when-to-use-which/}
\item How to make a User Flow diagram \url{https://www.lucidchart.com/blog/how-to-make-a-user-flow-diagram}
\end{itemize}

\textbf{Quote:}

'Each use case is represented as a sequence of simple steps, beginning with a user's goal and ending when that goal is fulfilled.' Usability.gov

